\subsection*{Every quartic abelian input with zero compact coordinate}
\begin{statement}{E.108}{Proved}{The complete peripheral quartic class-two image and explicit central coefficient maps}
Use the plus quartic, the actual section $s_R$, its class-two
lift $S_R$, and all tangent, coset and coefficient maps of
E.103 and E.105.  Here $R=\mathbb Z_\ell$ for an odd prime,
the product of these rings, or $R=\mathbb Q_\ell$ for any prime.
Put $F=\mathbb Q(\zeta_8)$, $H=G_F$, $V=T_RE_+$ and
$A=W\oplus V^3$ using exactly (E103.16).
The following criterion evaluates every abelian input
whose three compact coordinates are zero.  Thus its input
space, with the actual forward and inverse maps of E.103, is
\[
 \mathcal W_R=H^1(\mathbb Q,W)
   = (\widehat{F^\times}_R)^2\times
       \ker\bigl(N_{F/\mathbb Q}:\widehat{F^\times}_R
                                      \longrightarrow\widehat{\mathbb Q^\times}_R\bigr).
 \tag{E108.1}
\]
For rational $R$ the hats in this display mean the corresponding
power completions followed by scalar extension, with the bounded
continuous-cochain topology specified below.

Let $b:G_{\mathbb Q}\to W$ represent an input.  On $H$ write
$b=(b_i)_{i=1}^{11}$ in the peripheral order (E105.11), and
put $k=\kappa_2$.  Retain the literal matrix $N$ of (E105.12)
and vector $v$ of (E105.16).  For every $i<j$ define
\[
 p_{ij}=v_i-N_{ij}/8,\qquad q_{ij}=v_j-N_{ij}/8,\qquad
 X_{ij}=b_i+p_{ij}k,\qquad Y_{ij}=b_j+q_{ij}k.
 \tag{E108.2}
\]
All these divisions are defined in $R$.  The class of $b s_R$
lifts to the actual intrinsic class-two extension if and only if
all $55$ scalar classes $X_{ij}\cup Y_{ij}$ vanish in
$H^2_{\rm cts}(F,R(2))$.  The next formulas give a test on all
completed inputs, rather than leaving these cohomology classes
unevaluated.

Choose $\epsilon=1+\sqrt2$ and one generator $\pi_{\mathfrak p}$
of each prime ideal of $F$, with coherent roots.  A class in
$H^1(F,R(1))$ has the unique coordinates
\[
 x=x_\epsilon\kappa_\epsilon+
                 \sum_{\mathfrak p}x_{\mathfrak p}\kappa_{\pi_{\mathfrak p}},
 \qquad x_{\mathfrak p}\longrightarrow0.
 \tag{E108.3}
\]
For integral odd $R$ these are integral coordinates; for rational
$R$ they lie in $\mathbb Q_\ell$ and tend to zero there.  The
actual root-of-unity coboundary at two is retained in the proof.
For each pair $(X,Y)$ from (E108.2), denote the coordinates in
(E108.3) by $(a_\epsilon,a_{\mathfrak p})$ and
$(b_\epsilon,b_{\mathfrak p})$.
At $\mathfrak q\nmid\ell$ put $Q_{\mathfrak q}=N\mathfrak q$,
$e_{\mathfrak q}=v_\ell(Q_{\mathfrak q}-1)$, and choose a primitive
$\ell^{e_{\mathfrak q}}$th root $\eta_{\mathfrak q}$ in its
residue field, compatibly with the fixed roots already in $F$.
Thus at two its $2^{e_{\mathfrak q}-3}$th power is the residue
of the fixed $\zeta_8$.  For $u=\epsilon$ or $u=\pi_{\mathfrak p}$ with
$\mathfrak p\ne\mathfrak q$, define the bounded integer
$r_{\mathfrak q,u}$ by
\[
 \begin{gathered}
 \bar u^{(Q_{\mathfrak q}-1)/\ell^{e_{\mathfrak q}}}
       =\eta_{\mathfrak q}^{r_{\mathfrak q,u}},\qquad
       0\le r_{\mathfrak q,u}<\ell^{e_{\mathfrak q}},\\
 A_{\mathfrak q}=a_\epsilon r_{\mathfrak q,\epsilon}
        +\sum_{\mathfrak p\ne\mathfrak q}
                  a_{\mathfrak p}r_{\mathfrak q,\pi_{\mathfrak p}},\qquad
 B_{\mathfrak q}=b_\epsilon r_{\mathfrak q,\epsilon}
        +\sum_{\mathfrak p\ne\mathfrak q}
                  b_{\mathfrak p}r_{\mathfrak q,\pi_{\mathfrak p}},\\
 c_{\mathfrak q}(X,Y)=
       -\frac{A_{\mathfrak q}b_{\mathfrak q}
                      -B_{\mathfrak q}a_{\mathfrak q}}
                    {Q_{\mathfrak q}-1}.
 \end{gathered}
 \tag{E108.4}
\]
When $e_{\mathfrak q}=0$ all the indicated residues are zero.
The sums in this formula converge uniformly in the coefficient
congruence topology.  The exact lifting test is
\[
 \begin{array}{ll}
 R=\mathbb Z_\ell,\ \ell\ne2:
   &c_{\mathfrak q}(X_{ij},Y_{ij})\in\mathbb Z_\ell
       \text{ for every }\mathfrak q\nmid\ell,
       \quad c_{\mathfrak q}(X_{ij},Y_{ij})\longrightarrow0;\\[2pt]
 R=\mathbb Q_\ell:
   &c_{\mathfrak q}(X_{ij},Y_{ij})\longrightarrow0
                                      \text{ in }\mathbb Q_\ell.
 \end{array}
 \tag{E108.5}
\]
Apply these tests to all $i<j$, and componentwise for the product
of odd coefficient rings.  In the rational case convergence
includes one common bound on all denominators.  No closure is
taken in the support quotient used to derive (E108.5).
For an input consisting of elements of $F^\times$, let $u_i$
be the eleven elements in (E108.9).  The integral odd test
can equivalently be performed with ordinary finite tame symbols:
set $U_{ij}=u_i^8 2^{8v_i-N_{ij}}$ and
$V_{ij}=u_j^8 2^{8v_j-N_{ij}}$, whose exponents are integers,
and require the $\ell$-primary component of
$(-1)^{v_{\mathfrak q}(U_{ij})v_{\mathfrak q}(V_{ij})}
 U_{ij}^{v_{\mathfrak q}(V_{ij})}/
 V_{ij}^{v_{\mathfrak q}(U_{ij})}$ in the residue field to
be trivial at every prime ideal.  Only the union of the
numerator and denominator supports is needed.

Every input in (E108.1) represented by classes unramified outside
two passes this test.  This is an explicit $R^5$ family.  Every
input represented by finite-support elements of $F^\times$ passes
the rational-coefficient test.  The integral test has nonzero
obstructions even over the zero toric minimum.  In the input
coordinates of (E108.1), take
\[
       (a,b,c)=\left(1,1,\frac{3+\sqrt2}{3-\sqrt2}\right).
 \tag{E108.6}
\]
Its norm-one third coordinate has compact and toric-minimum
coordinates zero.  Its $(C_0,C_1)$ obstruction at the prime
$(3+\sqrt2)$ of norm $49$ is nonzero for $R=\mathbb Z_3$.
It therefore has no intrinsic integral three-primary
class-two lift.

Every nonempty fibre over a fixed input is the entire torsor
under $H^1_{\rm cts}(\mathbb Q,\Lambda^2 A)$, with the actual
central-cocycle action and inverse ratio of E.105.  The following
direct summands and their coordinate maps are evaluated here:
\[
 \begin{gathered}
 H^1(\mathbb Q,\Lambda^2 W)=
 \begin{cases}
   \mathbb Z_3^{30}\oplus(\mathbb Z/3)^{10},&R=\mathbb Z_3,\\
   \mathbb Z_\ell^{30},&R=\mathbb Z_\ell,\quad\ell\ne2,3,\\
   \mathbb Q_\ell^{30},&R=\mathbb Q_\ell;
 \end{cases}\\
 H^1(\mathbb Q,\Lambda^2 A)\simeq
 H^1(\mathbb Q,\Lambda^2 W)
 \oplus\mathcal R_{\mathcal P_{\mathbb Q}}(R)^6
 \oplus\mathcal R_{\{p:p\equiv1\ (4)\}}(R)^3\\
 \hspace{12mm}\oplus H^1(\mathbb Q,V(1))^{15}
 \oplus H^1(\mathbb Q,V(1)\chi_2)^{18}
 \oplus H^1(\mathbb Q,U_{\rm CM})^3.
 \end{gathered}
 \tag{E108.7}
\]
Here $\chi_2$ is the character of $\mathbb Q(\sqrt2)$,
$\mathcal R$ is the restricted coefficient module (E103.5)
or its rational $c_0$ version, and $U_{\rm CM}$ is the actual
rank-two coefficient kernel in (E108.24) below.
For the product of odd rings use the product of the first-line
answers, with its specified topology.
The remaining three types of CM cohomology factors in the last
line are retained with their actual coefficient projections and
sections.  No numerical evaluation of those factors is asserted.
The theorem computes the full zero-compact-coordinate image;
the mixed inputs and the compact arithmetic action retained in
E.23 and E.109 are still required for the complete quartic target.
\end{statement}

\begin{proof}
\emph{The actual input maps and root choices.}
Write $\Delta=\operatorname{Gal}(F/\mathbb Q)$, indexed by
$\sigma_k(\zeta_8)=\zeta_8^k$, $k=1,3,5,7$.
The twelve-cusp quotient and its inverse are (E103.14)--(E103.16),
including the actual compact map $j_*/4$.
For precision, if $\Lambda=R[\Delta](1)$ and
$I=\ker(\operatorname{aug}:\Lambda\to R(1))$, their toric maps are
\[
 \begin{gathered}
 W=\Lambda^3/\delta_{\rm diag}R(1),\qquad
 [v_1,v_2,v_3]\longmapsto
 (v_1-\delta t,v_2-\delta t,v_3-\delta t),\quad
 t=\operatorname{aug}(v_3)/4,\\
 \Lambda^2\oplus I\longrightarrow W,\qquad
 (a,b,c)\longmapsto[a,b,c].
 \end{gathered}
 \tag{E108.8}
\]
These are inverse equivariant coefficient maps.  Apply the
literal Shapiro cocycle (E103.20), then the projector
$1-\delta\operatorname{aug}/4$ to the third component.  Its
augmentation is the actual Kummer cocycle of the norm, so the
result represents exactly (E108.1).  Conversely restriction to
$H$ followed by the coordinate at $1\in\Delta$ recovers each
Kummer class, and the coboundary inverse in (E103.20) recovers
the original global cohomology class.

Under this map, the eleven restricted coordinate classes can
also be read directly.  With $k_i=1+2i$, they are
\[
 \begin{array}{ll}
 [b_{A_i}]=\kappa_{\sigma_{k_i}a/b},&0\le i\le3,\\
 {}[b_{B_i}]=\kappa_{\sigma_{k_i}b/b},&1\le i\le3,\\
 {}[b_{C_i}]=\kappa_{\sigma_{k_i}c/b},&0\le i\le3.
 \end{array}
 \tag{E108.9}
\]
The quotient sets $e_{B_0}=-\sum_{d\ne B_0}e_d$, which gives
the divisions by $b$ in every row.  The projector on the third
Shapiro component subtracts $\kappa_{Nc}/4$ from each of its
four coordinates; this is the retained coboundary because
$Nc=1$ in the completed group.  Thus (E108.9) is a formula for
cohomology classes.  Whenever its Kummer representatives replace
the cocycle produced by (E108.8), their difference is the literal
coboundary $dm$, determined uniquely by $R(1)^H=0$.  Equation
(E108.13) below transports the class-two coordinates by that
same $m$.  This supplies the representative comparison as well
as the class formula.

The class number and full units of $F$ were proved in E.103:
$\mathcal O_F^\times=\mu_8\times\langle\epsilon\rangle$ and
every prime ideal is principal.  Consequently
$F^\times=\mu_8\times\epsilon^{\mathbb Z}\times
\bigoplus_{\mathfrak p}\pi_{\mathfrak p}^{\mathbb Z}$.
Taking powers, then inverse limits, gives (E108.3) with the
uniform congruence topology.  Its inverse reads the valuations
and then the unit exponent.  For odd coefficients the coherent
roots of $\mu_8$ already belong to $F$ and contribute zero.
For rational two-primary coefficients choose the roots
$\zeta_{8\cdot2^n}$ of $\zeta_8$; then
\[
 \kappa_{\zeta_8^t}=d(t/8),\qquad
 2=\zeta_8^6\epsilon^2(1-\zeta_8)^4.
 \tag{E108.10}
\]
The second equality is checked in $\mathbb Q[X]/(X^4+1)$.
Thus the class of $k=\kappa_2$ has unit coordinate $2$ and
the coordinate $4$ at $\mathfrak p_2=(1-\zeta_8)$.
If the product roots on the right of (E108.10) differ from the
fixed E.105 roots of two by $\zeta_{2^n}^{\gamma_n}$, their
coherent exponents give $\gamma\in\mathbb Z_2$ and the exact
cocycle correction is $d(3/4+\gamma)$, with the sign determined
by taking the prescribed root divided by the product root.
This is removed by the displayed rational coboundary and its
actual BCH conjugation.  For odd coefficients the same ratio
has integral exponent and can be absorbed into the roots of
$\epsilon$, because two is invertible.
All preceding rational-two tangent and endpoint transports
remain exactly those of (E105.19); (E108.10) changes the
arithmetic unit basis, not the specified tangential origin.

\emph{All fifty-five arithmetic obstruction coordinates.}
On $H$, (E105.15)--(E105.20) give the actual action relative
to $S_R$ on the peripheral class-two group:
\[
 g(x,z)=(\chi_gx,\chi_g^2z+\chi_gk(g)L_v(x)),\qquad
 L_v(e_i)=e_i\wedge(N_{\bullet i}/8)+v\wedge e_i.
 \tag{E108.11}
\]
The symmetry of $N$ makes the $(i,j)$ row of $L_v$ equal to
$-q_{ij}$ on $e_i$, $p_{ij}$ on $e_j$, and zero on every
other basis vector.  Put $x=b_i$, $y=b_j$, $p=p_{ij}$,
$q=q_{ij}$.  The coefficient of (E105.23) is therefore
\[
 \begin{gathered}
 (\Omega_b)_{ij}=-q\,k\cup x+p\,k\cup y
                       +\tfrac12(x\cup y-y\cup x),\\
 (\Omega_b)_{ij}=(x+pk)\cup(y+qk)+d\Phi_{ij},\qquad
 \Phi_{ij}(g)=\tfrac12x(g)y(g)+q\,k(g)x(g)
                                      +\tfrac12pq\,k(g)^2.
 \end{gathered}
 \tag{E108.12}
\]
Here $x\cup y(g,h)=x(g)\chi_g y(h)$ and
$d\Phi(g,h)=\chi_g^2\Phi(h)-\Phi(gh)+\Phi(g)$.
To prove the identity, expand
$d(xy)=-(x\cup y+y\cup x)$,
$d(kx)=-(k\cup x+x\cup k)$ and
$d(k^2)=-2k\cup k$.  The two unwanted mixed terms and the
$k\cup k$ term cancel literally.  The certificate verifies
this polynomial identity and retains $p,q,pq/2$ for every
one of the $55$ pairs from the frozen E.105 matrix.

The representative change promised after (E108.9) is also
literal.  Conjugating a class-two cocycle $(b(g),z(g))$ by
$(-m,0)$ in the BCH group gives
\[
 \begin{split}
 b'(g)&=b(g)+(\chi_g-1)m,\\
 z'(g)&=z(g)+\chi_gk(g)L_v(m)
                         -\tfrac12(1+\chi_g)m\wedge b(g).
 \end{split}
 \tag{E108.13}
\]
Multiply the three factors
$(-m,0)(b,z)g(m,0)$ by (E105.14) to obtain this formula.
An additional central conjugator $w$ adds
$w-\chi_g^2w$.  Thus every replacement of roots, Shapiro
representatives or old coset coordinates has an explicit
coefficient map and its inverse.

For a prescribed global $b$, pull back the actual class-two
extension over the section $b s_R$.  Its kernel is the
abelian $G_{\mathbb Q}$-module $\Lambda^2A$.
On $H$ its obstruction is exactly (E108.12), and lies in
the direct summand $\Lambda^2 W$.  Necessity follows by
projecting a primitive to each of the $55$ coordinates.
Conversely, if $du_{ij}=X_{ij}\cup Y_{ij}$, the literal
central coordinate
\[
              z_{ij}=-u_{ij}-\Phi_{ij}
 \tag{E108.14}
\]
gives an $H$-section of this very pullback.  To descend it,
choose its four coset lifts $e_i$.  For a lift $e$ of $g$,
write $j=g\cdot i$, $h_i=g_j^{-1}gg_i$, and set
$B(e)=\sum_i e e_i s_H(h_i)^{-1}e_j^{-1}$ in its central
kernel.  Then $B(ef)=B(e)+gB(f)$ and $B(z)=4z$.
The formula $g\mapsto(-B(e)/4)e$ is a continuous section,
independent of $e$.  This is the actual four-coset formula
applied to the abelian-kernel pullback.  It leaves the fixed
abelian input and origin unchanged.  Equivalently,
restriction on its obstruction is injective because
corestriction followed by restriction on the source class
is multiplication by four, a unit in $R$.
This proves both directions of the $55$-cup criterion.

\emph{The finite base group over $F$ and its exact support limit.}
Guo--Qin--Song prove
$K_2(\mathcal O_F)=0$ in their Theorem 3.5
\cite[Theorem 3.5 and its proof]{E108_GuoQinSong}.
Their proof reduces the tame kernel to the three symbols
formed from $\zeta_8$ and $\zeta_8\epsilon$, then kills them
by Steinberg identities.  The identity
$(1-\zeta_8)(1+\zeta_8+\zeta_8^2)=1-\zeta_8^3$
is among the retained exact field checks.
Localizing at the primes above $\ell$ has cokernel contained
in the sum of their finite residue-unit groups, whose orders
are prime to $\ell$.  The additional prime above two has
residue-unit group trivial.  Hence, for
$S=\{\mathfrak p:\mathfrak p\mid2\ell\}\cup S_\infty$,
$K_2(\mathcal O_{F,S})/\ell^n=0$ for every $n$.

Kahn's proved Appendix 2, Proposition A2.1 identifies this
quotient with $H^2(\mathcal O_{F,S},\mu_{\ell^n}^{\otimes2})$
for every $\ell$, including two
\cite[pp.97--98, Proposition A2.1]{E108_Kahn1993}.
In particular the prime-power assertion uses his completed
injectivity argument, as well as the prime-coefficient case.
The arithmetic etale comparison and the absence of real
places give
\[
 H^2(G_{F,S},\mathbb Z/\ell^n(2))
   =H^3(G_{F,S},\mathbb Z/\ell^n(2))=0,
 \qquad H^1(G_{F,S},\mathbb Z/\ell^n(2))\text{ finite}.
 \tag{E108.15}
\]
These are the same finite localization, purity and Galois
comparison maps as in E.32 and E.105: Milne II.0.1(f),
II.1.7(b), II.2.9 and I.4.10(c), I.4.15
\cite{E_MilneADT}.
One may omit the prime above two from $S$ whenever
$\ell\ne2$, because its residue-unit group is trivial.

First take the filtered direct limit over finite sets of
removed primes, at fixed $\ell^n$.  The support term is
the direct sum of the complexes
$[\mathbb Z/\ell^n\xrightarrow{Q_{\mathfrak q}-1}
\mathbb Z/\ell^n]$ in degrees zero and one, with the
support shift by one.  The arithmetic Frobenius normalization
comes from $T(-1)$, where $T=\mathbb Z_\ell(2)$.
Both terms have surjective coefficient transitions.
Taking their derived inverse limit therefore gives the
actual complex $[\mathcal R_F\xrightarrow d\mathcal R_F]$,
where $\mathcal R_F$ consists of integral sequences tending
to zero, indexed by $\mathfrak q\nmid\ell$, and
$d(c)_{\mathfrak q}=(Q_{\mathfrak q}-1)c_{\mathfrak q}$.
The omitted prime above two, if present, contributes an
acyclic complex with differential one.
The finite degree-one base groups and (E108.15) kill the
adjacent derived-limit terms.  We obtain
\[
 H^2_{\rm cts}(F,\mathbb Z_\ell(2))
          \simeq\mathcal R_F/d\mathcal R_F.
 \tag{E108.16}
\]
The denominator is the image itself, with no closure.
This uses the exact inverse-limit maps of
\cite[Proposition 1.6 and Theorem 2.2]{E_JannsenContinuous}.
It retains possible classes invisible at every individual
finite coefficient level.

For rational coefficients the map
$C^r_{\rm cts}(G_F,\mathbb Z_\ell(2))\otimes\mathbb Q_\ell
\to C^r_{\rm cts}(G_F,\mathbb Q_\ell(2))$ is an isomorphism
in every degree: the image of a continuous cochain on the
compact space $G_F^r$ is bounded.  Multiplying by a single
power of $\ell$ makes it integral.  These maps commute with
differentials.  Thus rationalizing (E108.16) gives
\[
 H^2_{\rm cts}(F,\mathbb Q_\ell(2))
       \simeq c_0(\mathfrak q\nmid\ell,\mathbb Q_\ell)\big/
                d\,c_0(\mathfrak q\nmid\ell,\mathbb Q_\ell).
 \tag{E108.17}
\]
This proof includes $\ell=2$, since (E108.15) already
does.  It does not infer an unrestricted comparison by
passing to the limit of separate local vanishings.

\emph{The actual residue representative and its inverse.}
At each $\mathfrak q\nmid\ell$ choose tame inertia $\tau$
and arithmetic Frobenius $\phi$ with
$\phi\tau\phi^{-1}=\tau^{Q_{\mathfrak q}}$ and
$\kappa_{\pi_{\mathfrak q}}(\tau)=1$,
$\kappa_{\pi_{\mathfrak q}}(\phi)=0$.
The last condition is achieved by the corresponding inertia
translate of Frobenius.  A choice of prolongation sends the
fixed root of unity of order $\ell^{e_{\mathfrak q}}$ to
$\eta_{\mathfrak q}$; at two its image must respect the
already specified copy of $\mu_8$ in $F$.
For a pair $X,Y$, let $A'_{\mathfrak q},B'_{\mathfrak q}$
be its actual Frobenius values, and $a_{\mathfrak q},
b_{\mathfrak q}$ its inertia values.  The alternating BCH
cup has the central-zero splitting on inertia.  Direct
multiplication gives its tame relation defect with central
inertia correction $c$ as
\[
 Q_{\mathfrak q}\bigl((Q_{\mathfrak q}-1)c
       +A'_{\mathfrak q}b_{\mathfrak q}
       -B'_{\mathfrak q}a_{\mathfrak q}\bigr).
 \tag{E108.18}
\]
The support map divides this value by $Q_{\mathfrak q}$,
as it compares the value at $\tau^{Q_{\mathfrak q}}$ with
the value at $\tau$.  This is the typed cocycle representative
of the cup in (E108.16).

Finite-field Kummer theory gives
$\kappa_u(\phi)\equiv r_{\mathfrak q,u}
\pmod{Q_{\mathfrak q}-1}$ for every displayed unit $u$.
Write the difference as $(Q_{\mathfrak q}-1)t_{\mathfrak q,u}$
with $t_{\mathfrak q,u}\in\mathbb Z_\ell$.
The sums of these differences with the coordinates in
(E108.3) converge and are uniformly bounded.  Replacing
actual Frobenius values by $A_{\mathfrak q},B_{\mathfrak q}$
therefore changes the support representative by $d$ of a
restricted family: every term contains
$a_{\mathfrak q}$ or $b_{\mathfrak q}$, both tending to zero.
Changing $X,Y$ by $du,dw$ adds
$(Q_{\mathfrak q}-1)(u b_{\mathfrak q}-w a_{\mathfrak q})$.
Replacing $\phi$ by $\tau^r\phi$ adds
$r a_{\mathfrak q},r b_{\mathfrak q}$ to the two values
and leaves the alternating determinant unchanged.
Changing a compatible residue generator multiplies the
coordinate by a unit, together with the bounded root change
just described.  These are the actual maps proving
independence of all choices in (E108.4)--(E108.5).

The unique possible correction of (E108.18) is precisely
$c_{\mathfrak q}(X,Y)$ in (E108.4).  Membership in the
integral or rational support module is exactly (E108.5).
This proves the stated complete criterion.  It also gives
an inverse construction of a primitive: at each finite
coefficient level use these finite-support inertia values
in the localization complex.  Exactness and (E108.15)
give a global primitive, and the residue map in degree
one adjusts it to the prescribed inertia values.  The
remaining choices, after this adjustment, are torsors
under finite unramified degree-one cohomology groups.
Liftable finite systems have the finite intersection
property; compactness supplies compatible primitives.
One can select them by ordering the finite sets and
choosing the first extendable branch.  Their uniform
limit is the required continuous cochain $u_{ij}$.
For rational coefficients first multiply by one common
power of $\ell$, perform this construction, and divide
back.  Equations (E108.14) and the four-coset formula
then give the actual global section over the input.
For the product of odd rings the construction is
componentwise in the open-ideal topology of (E103.5).

For finite-support elements of $F^\times$ only finitely
many inertia coefficients occur, so (E108.5) always holds
rationally.  The two field elements $U_{ij},V_{ij}$ in the
statement represent $8X_{ij},8Y_{ij}$.  Their tame symbol
is exactly the residue determinant in (E108.18), after
multiplication by $64$; its sign has order at most two.
For odd coefficients $64$ is a unit and that sign vanishes.
This proves the finite-symbol version with its actual map.
For inputs unramified outside two, all the
shifted cups inflate from the base in (E108.15), hence
vanish integrally as well.  In the $S=\{2\}$ Kummer
coordinates each copy of $\widehat{F^\times}_R$ is
$R\kappa_\epsilon\oplus R\kappa_{1-\zeta_8}$.
The norm kills the first coordinate and maps the second
isomorphically to $R\kappa_2$.  The two unrestricted copies
and the norm kernel therefore give the asserted $R^5$.

For (E108.6), put $u=(3+\sqrt2)/(3-\sqrt2)$.
Its four conjugates are $u,u^{-1},u^{-1},u$, so $Nu=1$.
Equation (E108.9) gives zero on all $A$ and $B$ coordinates,
and $(\kappa_u,-\kappa_u,-\kappa_u,\kappa_u)$ on the
four $C$ coordinates.  Formula (E108.12) at $(C_0,C_1)$
has cohomology class
\[
                         -\tfrac14\kappa_u\cup\kappa_2.
 \tag{E108.19}
\]
The algebraic integer $3+\sqrt2$ has norm $49$; the quotient
of $\mathcal O_F$ by it is the field of $49$ elements,
since it lies over a degree-two prime above seven.
At this prime $u$ has valuation one, while two is a unit
of residue order three.  Its three-primary logarithm is
nonzero.  The support coordinate of (E108.19) is one
quarter of that logarithm and is nonzero modulo three.
The certificate retains the actual finite-field model,
the image of $\zeta_8$, the primitive residue root, and
this value.  The toric minimum is zero by (E103.4), proving
the claimed global nonlifting example with its exact map.

\emph{Peripheral central cohomology with actual character maps.}
Let $\chi_{-1},\chi_2,\chi_{-2}$ be the quadratic characters
of $\mathbb Q(i)$, $\mathbb Q(\sqrt2)$ and
$\mathbb Q(\sqrt{-2})$.  Their tables on $1,3,5,7$ are
\[
 \begin{array}{c|rrrr|r}
       &1&3&5&7&\text{multiplicity in }W(-1)\\ \hline
 1&1&1&1&1&2\\
 \chi_{-1}&1&-1&1&-1&3\\
 \chi_2&1&-1&-1&1&3\\
 \chi_{-2}&1&1&-1&-1&3
 \end{array}
 \tag{E108.20}
\]
For an orbit $D=A,B,C$ set
$u_{D,\chi}=\sum_{i=0}^3\chi(1+2i)e_{D_i}$ in the
twelve-cusp quotient.  Use $u_{A,1},u_{B,1}$ for the
trivial character and all three orbit columns for each
nontrivial character.  The omitted trivial column satisfies
$u_{C,1}=-u_{A,1}-u_{B,1}$.  Fourier inversion with its
factor $1/4$ proves these are a basis over every stated $R$.
The program retains its $11$ by $11$ matrix $P$ and
inverse, together with $\Lambda^2P$ and its inverse.
It checks them against all four actual permutation-quotient
matrices of E.105.  Their exterior character multiplicities
are
\[
           (1,\chi_{-1},\chi_2,\chi_{-2})
                         \quad\longmapsto\quad(10,15,15,15).
 \tag{E108.21}
\]
Thus these are explicit projections and sections of the
coefficient decomposition, including every integral lattice.

For odd $\ell$, restrict a scalar twisted Tate module
$\mathbb Z_\ell(2)\chi$ to $F$.  Equation (E108.15) and
the exact $\Delta$-invariants functor show that its finite
degree-two cohomology over $\mathbb Q$ unramified outside
$\{2,\ell\}$ is zero.  Global Euler characteristic then
gives one free degree-one coordinate for an odd character,
and zero free coordinates for an even character.  The
integral torsion is
$H^0(\mathbb Q,\mathbb Q_\ell/\mathbb Z_\ell(2)\chi)$.
It is $\mathbb Z/3$ exactly when $\ell=3$ and $\chi=1$,
and is zero in the other cases in (E108.20).
Indeed the cyclotomic square modulo $\ell$ can have order
at most two only for $\ell=3,5$; at five its nontrivial
quadratic character is ramified at five and equals none
of the three characters of conductor dividing eight.
At three the square character is trivial modulo three,
and its image on $1+3\mathbb Z_3$ has exactly one initial
power of three.  The asserted invariant group follows.

There is no further ramification in an integral degree-one
class: on tame inertia at $\mathfrak q\nmid\ell$ after
restriction to $F$, its value is killed by
$Q_{\mathfrak q}-1$ in a torsion-free lattice, so is zero.
Descent from $F$ adds no odd-primary inertia cohomology.
Thus the calculation is for unrestricted $G_{\mathbb Q}$.
For rational coefficients, restriction to $F$, the same
finite base calculation and scalar-extension argument
give the identical parity ranks also at two.  Equivalently
one can read these ranks through an explicit quadratic
restriction: Hain--Matsumoto Theorem 7.2 gives
$\dim H^1(G_{K,\{2\}},\mathbb Q_2(2))=r_2(K)$ for
$K=\mathbb Q,\mathbb Q(i),\mathbb Q(\sqrt2),
\mathbb Q(\sqrt{-2})$ \cite[Theorem 7.2]{E_HainMatsumoto}.
The restriction/corestriction projectors $(1\pm c)/2$
identify the nontrivial quadratic component with the
corresponding twisted cohomology over $\mathbb Q$.
The resulting ranks are respectively $0,1,0,1$, as above;
the tame argument already proved that no new unrestricted
degree-one class is added.

Here are actual inverse coordinates for these free factors.
For each odd $\chi$, choose a nonzero finite degree-one
class modulo $\ell$ and lift it successively, using the
vanishing finite degree-two groups.  The finite degree-one
groups are cyclic of order $\ell^n$, so the compatible
classes and their continuous cocycles give a generator
$u_\chi$.  A cocycle is recovered from its unique compatible
scalar coefficients modulo $\ell^n$, up to its coboundary.
At three the trivial-character torsion column is literally
$g\mapsto(\chi_3(g)^2-1)/3$.
For rational two-primary coefficients choose a nonzero
vector in each one-dimensional rational odd-character
space obtained by the preceding restriction calculation;
scalar coefficient and cocycle multiplication are inverse
linear maps.  These choices, followed by the stored
$\Lambda^2P$, give the first line of (E108.7), with its
forward and inverse topological maps.

\emph{The further evaluated central factors and the retained CM kernels.}
The six determinant projections and sections of (E105.8)
give six actual copies of $R(1)$.  Their degree-one
cohomology is $\mathcal R_{\mathcal P_{\mathbb Q}}(R)$,
with inverse coordinates the rational-prime valuations.
For the other factors use the actual CM automorphism
\[
            [i]:E_+\longrightarrow E_+,
                         \qquad(x,y)\longmapsto(-x,iy).
 \tag{E108.22}
\]
It satisfies $I^2=-1$ on $V$,
$gI=\chi_{-1}(g)Ig$, and preserves the Weil pairing.
Consequently $I:V\chi_{-1}\to V$ and its inverse $-I$
are explicit equivariant isomorphisms.  Applying them to
the character columns in (E108.20) gives
\[
                  W\otimes V^3
                          \simeq V(1)^{15}\oplus
                                      (V(1)\chi_2)^{18}.
 \tag{E108.23}
\]
No arithmetic action on a compact free-group generator
is replaced by this Tate-module computation.

In a symplectic basis $e,f$ of $V$, write
$I=\left(\begin{smallmatrix}a&b\\c&-a\end{smallmatrix}\right)$;
then $a^2+bc=-1$.  The symmetric bilinear contraction
$\ell(vw)=\langle Iv,w\rangle$ has target
$R(1)\chi_{-1}$.  In the basis $e^2,ef,f^2$ its row,
section and kernel are
\[
 \begin{gathered}
 \ell=(-c,a,b),\qquad
 s=(b/2,-a,-c/2)^{\rm t},\qquad \ell s=1,\\
 s\ell=(1+\operatorname{Sym}^2I)/2,\qquad
 U_{\rm CM}=\ker\ell=(1-s\ell)\operatorname{Sym}^2V,\\
 \operatorname{Sym}^2V\longrightarrow
         R(1)\chi_{-1}\oplus U_{\rm CM},\quad
 z\longmapsto(\ell z,z-s\ell z),\qquad
 (r,u)\longmapsto sr+u.
 \end{gathered}
 \tag{E108.24}
\]
The certificate proves these polynomial identities modulo
$a^2+bc+1$ and retains all maps.  Their equivariance follows
from the displayed CM relation and the Weil pairing, so
they hold over the actual arithmetic source.

Shapiro for $\mathbb Q(i)/\mathbb Q$, followed by the
anti-invariant projector $(1-c)/2$, identifies
$H^1(\mathbb Q,R(1)\chi_{-1})$ with the completed norm
kernel of $\mathbb Q(i)^\times$.  The Gaussian integer
ring is principal and its units have order four.
For a prime $p\equiv1\pmod4$, choose a factor $\varpi_p$;
the class $\varpi_p/\bar\varpi_p$ has valuation pair
$(1,-1)$.  These are the actual basis columns and their
inverse reads the valuation at $\varpi_p$.  At an inert or
ramified prime the unique prime is fixed by conjugation
and contributes zero to the anti-invariant part.
This proves the three restricted factors in (E108.7).
For completeness the principal-ring assertion follows
from the Euclidean norm on $\mathbb Z[i]$: choose the
nearest Gaussian integer to the real and imaginary parts
of a quotient, giving remainder norm at most half the
divisor norm.  Its units and the stated split-prime
factorization follow directly.  The rational-two root
torsion disappears through its explicit Kummer coboundary.

Finally $A^{G_{\mathbb Q}}=0$ by E.105.  The abelian
conjugator aligning two sections over the same class is
therefore unique.  Lift it using the fixed BCH section;
then the ratio of the two class-two sections is a central
cocycle, and central conjugation is exactly its ordinary
coboundary.  This gives all fibres, with the literal
inverse coordinate map and the continuous topology proved
at the end of E.105.  Applying the direct coefficient maps
(E108.20)--(E108.24) proves the entire decomposition
(E108.7).  The unevaluated CM groups in that display and
the actual compact and mixed arithmetic obstruction are
retained explicitly; neither the existence of the
$R^5$ family nor the $55$-coordinate test assigns values
to them.
\end{proof}
